\section{Discussion}
\label{sec:discussion}

\paragraph{Two simultaneous conclusions.}
Candidate B remains negative at the raw level and positive at the grouped
contract level.  The split implementation exposes finer singleton deletions
than the bundled implementation, so raw repair non-canonicity remains.  Once
the artifact-defined reportability contract, residual faithfulness, and group
soundness are fixed, the grouped report agrees with the bundled contract repair
family.

\paragraph{Contract scope.}
An artifact-defined contract takes recorded bundled residuals as the reference
predicate.  A semantic contract would independently justify those residuals as
the intended social-choice problems; a Lean-level contract would tie that
meaning to formal definitions and theorems.  M3 proves reportability relative
to the chosen predicate and does not silently perform either upgrade.

\paragraph{Why reportability must be earned.}
Implementation labels are not automatically explanation categories.
Atomicity earns reportability syntactically; group soundness earns it for
non-atomic interfaces; deletion monotonicity shows when group soundness is also
necessary for grouped correctness.  The finite examples in the Lean core show
that these distinctions are substantive.

\paragraph{Program boundary.}
The paper fits a larger program moving from auditable finite atlases toward
broader social-choice claims.  It contributes the repair/report component, not
a family-scale bridge theorem.  New theorem families require their own
reference residuals, evidence packages, and transfer arguments before M3 can
consume them.

\paragraph{Conclusion placeholder.}
M1.5 shows why reportability cannot be assumed; M3 shows how reportability can
be earned.
